\begin{theorem}
For any \(s\ge 3\), there is a positive constant \(c_s\) such that for any
\(k\ge 2\),
\[
  r(s,k)\ge c_s\frac{k^{s-1}}{(\log k)^{2s-4}}.
\]
\end{theorem}
\begin{proof}
Fix \(s\ge3\) and put \(t=s-1\).  Then \(t\ge2\).  We first prove the
claimed lower bound for all sufficiently large \(k\), allowing constants to
depend on this fixed \(s\).

Let \(C=C(t)\) be the constant from
\(\noderef{StarProductForwardIndependentBound}\), enlarged if necessary.
For every sufficiently large \(k\), choose a power \(q\) of \(2\) such that
\[
  Cq(\log q)^2\le k\le 4Cq(\log q)^2. \tag{1}
\]
Such a \(q\) exists because consecutive powers of \(2\) differ by a bounded
factor and \((\log(2q)/\log q)^2\to1\).  Let \(K\) be the finite field of
order \(q\), let \(G=G(t,K)\), and let \(D^*(G)\) be the star product digraph
from \(\noderef{StarProductDigraph}\).

By \(\noderef{StarProductDigraphTransitiveFree}\), \(D^*(G)\) is
\(T_{t+1}\)-free.  Since \(t+1=s\), it is \(T_s\)-free.  By
\(\noderef{StarProductForwardIndependentBound}\) and the lower bound in (1),
\[
  \fwi_k(D^*(G))\le (Cq^t)^k. \tag{2}
\]

Apply \(\noderef{DigraphToGraphIndependentSetBound}\) to \(D^*(G)\).  We
obtain a \(K_s\)-free graph \(\Gamma^*\) on the same vertex set and
\[
  i_k(\Gamma^*)
  \le
  \left(\frac e k\right)^k \fwi_k(D^*(G))
  \le
  \left(\frac e k\right)^k(Cq^t)^k. \tag{3}
\]

We next estimate the number of vertices.  By
\(\noderef{PolarityGraphParameters}\), the polarity graph \(G(t,K)\) has
\[
  |V(G)|=\frac{q^{t+1}-1}{q-1},\qquad
  d=\frac{q^t-1}{q-1}.
\]
For all sufficiently large \(q\), these satisfy
\[
  |V(G)|\ge q^t/2,\qquad d\ge q^{t-1}/2.
\]
The vertices of \(D^*(G)\) are ordered edges of \(G\).  Therefore
\(\noderef{ProductDigraphVertexCard}\) gives
\[
  |V(D^*(G))|=|V(G)|d\ge q^{2t-1}/4. \tag{4}
\]

Set
\[
  p=\frac{k}{eCq^t}.
\]
By the upper bound in (1) and \(t\ge2\), increasing the large-\(k\) threshold
if necessary gives \(0\le p\le1\).  From (3),
\[
  p^k i_k(\Gamma^*)\le1.
\]
Thus \(\noderef{SamplingKsFreeRamseyBound}\) applies and yields
\[
  r(s,k)> p|V(D^*(G))|-1.
\]
Using (4), we get
\[
  r(s,k)
  >
  \frac{k}{eCq^t}\cdot\frac{q^{2t-1}}4-1
  \ge
  \frac{kq^{t-1}}{20C}-1, \tag{5}
\]
again after enlarging the threshold only by a constant depending on \(s\).

It remains to compare \(q\) with \(k\).  From the upper bound in (1) and
\(\log q\le \log k\) for large \(k\), there is a constant \(c'_s>0\) such
that
\[
  q\ge c'_s\frac{k}{(\log k)^2}. \tag{6}
\]
Substituting (6) into (5) gives a constant \(c''_s>0\) such that, for all
sufficiently large \(k\),
\[
  r(s,k)
  \ge
  c''_s\frac{k^t}{(\log k)^{2t-2}}.
\]
Since \(t=s-1\), this is
\[
  r(s,k)
  \ge
  c''_s\frac{k^{s-1}}{(\log k)^{2s-4}}. \tag{7}
\]
The subtraction of \(1\) in (5) is absorbed into the positive main term by
decreasing \(c''_s\) and increasing the large-\(k\) threshold.

Finally, handle the finitely many \(k\) below that threshold.  By
\(\noderef{RamseyNumber}\), \(r(s,k)\) is a positive integer, and
\(\log k>0\) for \(k\ge2\).  Hence the finitely many numbers
\[
  r(s,k)\frac{(\log k)^{2s-4}}{k^{s-1}}
\]
with \(2\le k\) below the threshold have a positive minimum, if there are any.
Choose \(c_s>0\) no larger than this minimum and no larger than the constant
in (7).  Then (7) proves the desired inequality for large \(k\), and the
choice of \(c_s\) proves it for the remaining \(k\ge2\).
\end{proof}
